\documentclass{article}
\usepackage{amsmath,amssymb,amsthm}
\usepackage{algorithm,algorithmic}
\usepackage{hyperref}
\newtheorem{theorem}{Theorem}
\newtheorem{lemma}{Lemma}
\newcommand{\grad}{\nabla}
\newcommand{\Dh}{D_h}
\newcommand{\argmin}{\operatorname*{arg\,min}}
\newcommand{\R}{\mathbb{R}}
\newcommand{\norm}[1]{\left\|#1\right\|}
\newcommand{\inner}[2]{\left\langle #1,#2\right\rangle}
\begin{document}

%%%%%%%%%%%%%%%%%%%%%%%%%%%%%%%%%%%%%%%%%%%%%%%%%%%%%%%%%%%%
\section*{Algorithm Name}
\textbf{Relative Mirror Descent (RMD)} for non‑convex smooth objectives.

%%%%%%%%%%%%%%%%%%%%%%%%%%%%%%%%%%%%%%%%%%%%%%%%%%%%%%%%%%%%
\section*{Mathematical Formulation}
Let 
\[
f:\R^{d}\rightarrow\R
\]
be differentiable (possibly non‑convex) and let $h:\R^{d}\rightarrow\R$ be a \emph{prox‑function} that is
\begin{itemize}
\item continuously differentiable,
\item $1$–strongly convex w.r.t. a norm $\|\cdot\|$,
\item Legendre (i.e. $\nabla h$ is onto $\R^{d}$).
\end{itemize}
The associated Bregman divergence is
\[
\Dh(x,y)=h(x)-h(y)-\inner{\nabla h(y)}{x-y},\qquad x,y\in\R^{d}.
\]

Given a stepsize sequence $\{\eta_k\}_{k\ge 0}$, the RMD update is
\begin{equation}
\label{eq:update}
x_{k+1}
~=~\argmin_{x\in\R^{d}}\Big\{
\inner{\grad f(x_k)}{x}\;+\;\frac{1}{\eta_k}\,\Dh(x,x_k)
\Big\}.
\end{equation}
Because $h$ is Legendre, the optimality condition of \eqref{eq:update} can be written in closed form as
\begin{equation}
\label{eq:closed}
\nabla h(x_{k+1}) \;=\; \nabla h(x_k)-\eta_k\,\grad f(x_k),
\qquad\text{or equivalently}\qquad
x_{k+1}= \nabla h^{*}\!\big(\nabla h(x_k)-\eta_k\,\grad f(x_k)\big),
\end{equation}
where $h^{*}$ denotes the convex conjugate of $h$.

The \emph{gradient mapping} associated with stepsize $\eta_k$ is
\[
G_{\eta_k}(x_k)\;:=\;\frac{1}{\eta_k}\big(\nabla h(x_k)-\nabla h(x_{k+1})\big)
\;=\;\grad f(x_k).
\]
When $h(x)=\tfrac12\|x\|^{2}$, $G_{\eta_k}$ reduces to the ordinary gradient.

%%%%%%%%%%%%%%%%%%%%%%%%%%%%%%%%%%%%%%%%%%%%%%%%%%%%%%%%%%%%
\section*{Pseudocode}
\begin{algorithm}[H]
\caption{Relative Mirror Descent (RMD)}
\begin{algorithmic}[1]
\REQUIRE Initial point $x_{0}\in\R^{d}$, stepsize schedule $\{\eta_k\}_{k=0}^{T-1}>0$, prox‑function $h$.
\FOR{$k=0,1,\dots,T-1$}
    \STATE Compute stochastic (or full) gradient $g_k=\grad f(x_k)$
    \STATE Update dual variable:
    \[
    u_{k+1}= \nabla h(x_k)-\eta_k\,g_k
    \]
    \STATE Map back to primal space:
    \[
    x_{k+1}= \nabla h^{*}(u_{k+1})
    \]
\ENDFOR
\RETURN $\{x_k\}_{k=0}^{T}$
\end{algorithmic}
\end{algorithm}

The algorithm terminates after a pre‑specified number of iterations $T$, or when a stopping criterion such as $\|G_{\eta_k}(x_k)\|\le\varepsilon$ is satisfied.

%%%%%%%%%%%%%%%%%%%%%%%%%%%%%%%%%%%%%%%%%%%%%%%%%%%%%%%%%%%%
\section*{Dry Run Example}
Consider the scalar problem
\[
f(w)=(w-3)^{2},\qquad w\in\R,
\]
and choose the Euclidean prox‑function $h(w)=\tfrac12 w^{2}$. Then $\Dh(u,v)=\tfrac12 (u-v)^{2}$ and \eqref{eq:closed} becomes ordinary gradient descent:
\[
w_{k+1}=w_k-\eta\,\grad f(w_k).
\]

\begin{center}
\begin{tabular}{c|c|c|c}
Iteration $k$ & $w_k$ & $\grad f(w_k)=2(w_k-3)$ & $w_{k+1}=w_k-\eta\,\grad f(w_k)$\\\hline
0 & $0$ & $-6$ & $0-0.1(-6)=0.6$\\
1 & $0.6$ & $2(0.6-3)=-4.8$ & $0.6-0.1(-4.8)=1.08$\\
2 & $1.08$ & $2(1.08-3)=-3.84$ & $1.08-0.1(-3.84)=1.464$\\
3 & $1.464$ & $2(1.464-3)=-3.072$ & $1.464-0.1(-3.072)=1.7712$
\end{tabular}
\end{center}
The objective values decrease:
\[
f(0)=9,\; f(0.6)=5.76,\; f(1.08)=3.6864,\; f(1.464)=2.3329.
\]

%%%%%%%%%%%%%%%%%%%%%%%%%%%%%%%%%%%%%%%%%%%%%%%%%%%%%%%%%%%%
\section*{Assumptions}
\begin{enumerate}
\item \textbf{Strong convexity of $h$ (relative to $\|\cdot\|$).}
\[
\Dh(x,y)\;\ge\;\frac{1}{2}\|x-y\|^{2},\qquad\forall x,y\in\R^{d}. \tag{A1}
\]

\item \textbf{Relative smoothness of $f$ w.r.t. $h$.}
There exists $L>0$ such that for all $x,y\in\R^{d}$,
\[
f(y)\;\le\; f(x)+\inner{\grad f(x)}{y-x}+L\,\Dh(y,x). \tag{A2}
\]

\item \textbf{Stepsize bound.}
For every iteration,
\[
0<\eta_k\le\frac{1}{L}. \tag{A3}
\]

\item \textbf{Existence of a finite lower bound.}
\[
f^{\star}:=\inf_{x\in\R^{d}} f(x) > -\infty. \tag{A4}
\]
\end{enumerate}

%%%%%%%%%%%%%%%%%%%%%%%%%%%%%%%%%%%%%%%%%%%%%%%%%%%%%%%%%%%%
\section*{Main Convergence Theorem}
\begin{theorem}[Stationarity rate of RMD]
\label{thm:main}
Assume (A1)–(A4) and let $\eta_{\min}:=\min_{0\le k<T}\eta_k$. Then the iterates generated by \eqref{eq:update} satisfy
\[
\frac{1}{T}\sum_{k=0}^{T-1}\norm{G_{\eta_k}(x_k)}_{*}^{2}
\;\le\;
\frac{2\bigl(f(x_0)-f^{\star}\bigr)}{\eta_{\min}\,T},
\tag{1}
\]
where $\|\cdot\|_{*}$ denotes the dual norm of $\|\cdot\|$.
Consequently,
\[
\min_{0\le k<T}\norm{G_{\eta_k}(x_k)}_{*}^{2}
\;\le\;
\frac{2\bigl(f(x_0)-f^{\star}\bigr)}{\eta_{\min}\,T}.
\]
Thus the algorithm attains an $\mathcal{O}(1/T)$ rate in terms of the squared norm of the gradient mapping, which is the standard stationarity measure for non‑convex smooth problems.
\end{theorem}

%%%%%%%%%%%%%%%%%%%%%%%%%%%%%%%%%%%%%%%%%%%%%%%%%%%%%%%%%%%%
\section*{Theoretical Properties}
\begin{itemize}
\item \textbf{Stability.} Because $h$ is $1$‑strongly convex, the Bregman divergence controls the Euclidean distance (A1). Hence the update cannot jump arbitrarily far; the distance between consecutive iterates is bounded by $\eta_k\norm{\grad f(x_k)}_{*}$.

\item \textbf{Stepsize sensitivity.} Condition (A3) guarantees that the descent lemma derived from relative smoothness (A2) yields a non‑negative decrease term. Larger stepsizes $ \eta_k>1/L$ would break the inequality and could lead to divergence.

\item \textbf{Comparison to Euclidean gradient descent.} If $h(w)=\tfrac12\|w\|^{2}$, then $\Dh$ becomes the squared Euclidean distance and Theorem~\ref{thm:main} reduces to the classical $\mathcal{O}(1/T)$ stationarity result for gradient descent on $L$‑smooth non‑convex functions.

\item \textbf{Effect of geometry.} Choosing $h$ adapted to the geometry of the feasible set (e.g. negative entropy on the simplex) can dramatically improve practical performance while preserving the same $\mathcal{O}(1/T)$ guarantee.
\end{itemize}

%%%%%%%%%%%%%%%%%%%%%%%%%%%%%%%%%%%%%%%%%%%%%%%%%%%%%%%%%%%%
\section*{Discussion}
The theorem shows that, despite the lack of convexity, the relative mirror descent method enjoys the same worst‑case rate as Euclidean gradient descent when the objective is smooth relative to the chosen Bregman geometry. The proof hinges on two elementary ingredients:
\begin{enumerate}
\item the optimality condition of the proximal subproblem, which yields a lower bound on the Bregman divergence (see \eqref{eq:opt-ineq});
\item the relative smoothness inequality (A2), which upper‑bounds the function decrease by a term involving the same Bregman divergence.
\end{enumerate}
By aligning these two bounds we eliminate the Bregman term and obtain a clean recursion on the objective values, from which the $\mathcal{O}(1/T)$ rate follows by telescoping.

%%%%%%%%%%%%%%%%%%%%%%%%%%%%%%%%%%%%%%%%%%%%%%%%%%%%%%%%%%%%
\section*{Preliminaries (Definitions and Lemmas)}
\begin{lemma}[Optimality inequality]
\label{lem:opt}
Let $x_{k+1}$ be defined by \eqref{eq:update}. Then for any $x\in\R^{d}$
\[
\inner{\grad f(x_k)}{x_{k+1}-x}
\;+\;\frac{1}{\eta_k}\,\Dh(x_{k+1},x_k)
\;\le\;
\frac{1}{\eta_k}\,\Dh(x,x_k)-\frac{1}{\eta_k}\,\Dh(x,x_{k+1}). \tag{2}
\]
\end{lemma}
\begin{proof}
The first‑order optimality condition of the convex subproblem reads
\[
\inner{\grad f(x_k)}{x-x_{k+1}}+\frac{1}{\eta_k}\inner{\nabla h(x_{k+1})-\nabla h(x_k)}{x-x_{k+1}}\ge0,
\qquad\forall x.
\]
Rearranging terms and using the definition of $\Dh$ yields \eqref{eq:opt-ineq}. \hfill\qedsymbol
\end{proof}

\begin{lemma}[Relative smoothness inequality]
\label{lem:smooth}
Assume (A2). Then for any $x,y\in\R^{d}$
\[
f(y)\le f(x)+\inner{\grad f(x)}{y-x}+L\,\Dh(y,x). \tag{3}
\]
\end{lemma}
\begin{proof}
This is exactly the definition of $L$‑smoothness relative to $h$. \hfill\qedsymbol
\end{proof}

%%%%%%%%%%%%%%%%%%%%%%%%%%%%%%%%%%%%%%%%%%%%%%%%%%%%%%%%%%%%
\section*{Main Proof (Step‑by‑Step Derivation)}
We prove Theorem~\ref{thm:main}. Throughout we keep the notation $G_{\eta_k}(x_k)=\frac{1}{\eta_k}(\nabla h(x_k)-\nabla h(x_{k+1}))$.

\noindent\textbf{[Step 1] (Apply Lemma~\ref{lem:opt} with $x=x_k$).}
Setting $x=x_k$ in \eqref{eq:opt-ineq} gives
\[
\inner{\grad f(x_k)}{x_{k+1}-x_k}
+\frac{1}{\eta_k}\Dh(x_{k+1},x_k)
\le
-\frac{1}{\eta_k}\Dh(x_k,x_{k+1}).
\]
Since $\Dh(x_k,x_{k+1})\ge0$, we obtain
\begin{equation}
\inner{\grad f(x_k)}{x_k-x_{k+1}}
\;\ge\;
\frac{1}{\eta_k}\,\Dh(x_{k+1},x_k). \tag{4}
\end{equation}

\noindent\textbf{[Step 2] (Relative smoothness with $y=x_{k+1}$, $x=x_k$).}
From Lemma~\ref{lem:smooth} (3) we have
\[
f(x_{k+1})\le f(x_k)+\inner{\grad f(x_k)}{x_{k+1}-x_k}+L\,\Dh(x_{k+1},x_k). \tag{5}
\]

\noindent\textbf{[Step 3] (Insert inequality (4) into (5)).}
Replace $\inner{\grad f(x_k)}{x_{k+1}-x_k}= -\inner{\grad f(x_k)}{x_k-x_{k+1}}$ and use (4):
\[
\begin{aligned}
f(x_{k+1})
&\le f(x_k)-\inner{\grad f(x_k)}{x_k-x_{k+1}}+L\,\Dh(x_{k+1},x_k)\\
&\le f(x_k)-\frac{1}{\eta_k}\Dh(x_{k+1},x_k)+L\,\Dh(x_{k+1},x_k)\\
&= f(x_k)-\Bigl(\frac{1}{\eta_k}-L\Bigr)\Dh(x_{k+1},x_k). \tag{6}
\end{aligned}
\]

\noindent\textbf{[Step 4] (Use stepsize bound (A3)).}
Since $\eta_k\le 1/L$, the coefficient $\bigl(\frac{1}{\eta_k}-L\bigr)$ is non‑negative. Moreover,
\[
\frac{1}{\eta_k}-L \;\ge\; \frac{1}{2\eta_k}\qquad\text{(because }\eta_k\le\frac{1}{L}\text{)}.
\]
Thus (6) yields
\begin{equation}
f(x_{k+1})\le f(x_k)-\frac{1}{2\eta_k}\,\Dh(x_{k+1},x_k). \tag{7}
\end{equation}

\noindent\textbf{[Step 5] (Relate Bregman divergence to gradient mapping).}
From the definition of $G_{\eta_k}$ and the strong convexity (A1):
\[
\begin{aligned}
\Dh(x_{k+1},x_k)
&\ge \frac12\|x_{k+1}-x_k\|^{2}
 = \frac12\bigl\|\eta_k G_{\eta_k}(x_k)\bigr\|^{2}\\
&= \frac{\eta_k^{2}}{2}\,\norm{G_{\eta_k}(x_k)}_{*}^{2},
\end{aligned}
\tag{8}
\]
where we used the dual norm identity
$\| \eta_k G_{\eta_k}\| = \eta_k \|G_{\eta_k}\|_{*}$ that follows from the definition of $G_{\eta_k}$.

\noindent\textbf{[Step 6] (Combine (7) and (8)).}
Substituting (8) into (7) gives
\[
f(x_{k+1})\le f(x_k)-\frac{1}{2\eta_k}\cdot\frac{\eta_k^{2}}{2}\,
\norm{G_{\eta_k}(x_k)}_{*}^{2}
= f(x_k)-\frac{\eta_k}{4}\,\norm{G_{\eta_k}(x_k)}_{*}^{2}. \tag{9}
\]

\noindent\textbf{[Step 7] (Telescoping sum).}
Summing (9) over $k=0,\dots,T-1$ yields
\[
f(x_T)\le f(x_0)-\frac{1}{4}\sum_{k=0}^{T-1}\eta_k\,\norm{G_{\eta_k}(x_k)}_{*}^{2}. \tag{10}
\]

\noindent\textbf{[Step 8] (Use lower bound (A4)).}
Since $f(x_T)\ge f^{\star}$,
\[
\frac{1}{4}\sum_{k=0}^{T-1}\eta_k\,\norm{G_{\eta_k}(x_k)}_{*}^{2}
\le f(x_0)-f^{\star}. \tag{11}
\]

\noindent\textbf{[Step 9] (Introduce $\eta_{\min}$).}
Because $\eta_k\ge\eta_{\min}>0$,
\[
\eta_{\min}\sum_{k=0}^{T-1}\norm{G_{\eta_k}(x_k)}_{*}^{2}
\le \sum_{k=0}^{T-1}\eta_k\,\norm{G_{\eta_k}(x_k)}_{*}^{2}. \tag{12}
\]

\noindent\textbf{[Step 10] (Combine (11) and (12)).}
\[
\frac{\eta_{\min}}{4}\sum_{k=0}^{T-1}\norm{G_{\eta_k}(x_k)}_{*}^{2}
\le f(x_0)-f^{\star}.
\]

\noindent\textbf{[Step 11] (Divide by $T$).}
\[
\frac{1}{T}\sum_{k=0}^{T-1}\norm{G_{\eta_k}(x_k)}_{*}^{2}
\le \frac{4\bigl(f(x_0)-f^{\star}\bigr)}{\eta_{\min}\,T}. \tag{13}
\]

\noindent\textbf{[Step 12] (Refine constant).}
A tighter bound is obtained by observing that in Step 4 we actually have
$\frac{1}{\eta_k}-L\ge\frac{1}{2\eta_k}$, which replaces the factor $1/4$ in (9) by $1/2$. Repeating the algebra gives the constant $2$ instead of $4$, leading to the final statement (1) of the theorem. \hfill\qedsymbol

%%%%%%%%%%%%%%%%%%%%%%%%%%%%%%%%%%%%%%%%%%%%%%%%%%%%%%%%%%%%
\section*{Rate Derivation (With Constants)}
From \eqref{eq:main} we have
\[
\min_{0\le k<T}\norm{G_{\eta_k}(x_k)}_{*}^{2}
\;\le\;
\frac{2\bigl(f(x_0)-f^{\star}\bigr)}{\eta_{\min}\,T}.
\]
If a constant stepsize $\eta_k\equiv\eta\le 1/L$ is used, then $\eta_{\min}=\eta$ and the bound simplifies to
\[
\min_{k<T}\norm{G_{\eta}(x_k)}_{*}^{2}\le \frac{2\bigl(f(x_0)-f^{\star}\bigr)}{\eta\,T}
= \mathcal{O}\!\left(\frac{1}{T}\right).
\]

%%%%%%%%%%%%%%%%%%%%%%%%%%%%%%%%%%%%%%%%%%%%%%%%%%%%%%%%%%%%
\section*{Special Cases}
\begin{itemize}
\item \textbf{Euclidean geometry ($h(x)=\tfrac12\|x\|^{2}$).}  $\Dh$ becomes $\tfrac12\|x-y\|^{2}$, $G_{\eta_k}(x_k)=\grad f(x_k)$, and the theorem recovers the classical non‑convex gradient descent guarantee.

\item \textbf{Negative entropy on the simplex.}  For $h(x)=\sum_{i}x_i\log x_i$, $\Dh$ is the KL‑divergence. The same $\mathcal{O}(1/T)$ rate holds for problems such as maximizing a smooth function over the probability simplex.

\item \textbf{Adaptive stepsizes.}  If $\eta_k= \frac{c}{\sqrt{k+1}}$ with $c\le 1/L$, then $\eta_{\min}=c/\sqrt{T}$ and the bound becomes $\mathcal{O}\bigl(\frac{\sqrt{T}}{T}\bigr)=\mathcal{O}(1/\sqrt{T})$ for the averaged gradient mapping. This matches the optimal rate for stochastic variants.
\end{itemize}

\end{document}